\documentclass[10pt]{article}

\usepackage[utf8]{inputenc}
\usepackage[T1]{fontenc}
\usepackage{amsmath,amssymb,mathtools}
\usepackage{booktabs}
\usepackage{enumitem}
\usepackage[margin=0.78in]{geometry}
\usepackage{xcolor}
\usepackage[most]{tcolorbox}
\usepackage{hyperref}

\definecolor{BootcampBlue}{RGB}{31,78,121}
\definecolor{BootcampGold}{RGB}{215,155,35}
\hypersetup{colorlinks=true,linkcolor=BootcampBlue,urlcolor=BootcampBlue}
\setlist{itemsep=1pt,topsep=2pt,parsep=0pt,partopsep=0pt}
\setlength{\parindent}{0pt}
\setlength{\parskip}{3pt}
\setlength{\abovedisplayskip}{5pt plus 2pt minus 2pt}
\setlength{\belowdisplayskip}{5pt plus 2pt minus 2pt}

\newcommand{\R}{\mathbb{R}}
\newcommand{\E}{\mathbb{E}}
\newcommand{\Var}{\operatorname{Var}}
\newcommand{\Cov}{\operatorname{Cov}}
\newcommand{\trans}{^{\mathsf T}}
\newcommand{\norm}[1]{\left\lVert#1\right\rVert}
\newcommand{\argmin}{\operatorname*{arg\,min}}
\newcommand{\argmax}{\operatorname*{arg\,max}}
\newcommand{\proj}{\operatorname{proj}}
\newcommand{\diag}{\operatorname{diag}}

\newtcolorbox{deliverablebox}[1]{breakable,colback=BootcampGold!8,
  colframe=BootcampGold!80!black,boxrule=0.8pt,arc=1mm,
  title=\textbf{#1},fonttitle=\bfseries}
\newtcolorbox{conceptbox}[1]{breakable,colback=BootcampBlue!4,
  colframe=BootcampBlue,boxrule=0.8pt,arc=1mm,
  title=\textbf{#1},fonttitle=\bfseries}

\title{Summer Bootcamp -- Session 3\\Hands-On Optimization Exercises}
\author{Nathan Sauldubois}
\date{August 2026}

\begin{document}
\maketitle

\section*{Instructions}

Exercises 1--3 concern one-dimensional unconstrained optimization. Exercises
4--6 move to several variables. Exercises 7--9 introduce equality and
inequality constraints. Exercises 10--11 are algorithmic implementations.
Exercises 12--14 apply the theory to portfolio choice. The integrated Python
project is provided separately in the
\texttt{Portfolio Optimization Project} folder.

For analytical exercises, justify every candidate and distinguish necessary
from sufficient conditions. For Python work, use \texttt{NumPy},
\texttt{pandas}, \texttt{matplotlib}, and, where explicitly allowed,
\texttt{scipy.optimize}. Set all random seeds and report convergence criteria.
Do not use future test observations when estimating a portfolio.

\begin{conceptbox}{Notation and supplemental definitions}
The following notation completes the material explicitly introduced in the
session notes.
\begin{itemize}
  \item $\mathbf1=(1,\ldots,1)\trans\in\R^d$ is the vector of ones.
  For a symmetric positive-definite matrix $Q$, write
  $\lambda_{\min}(Q)$ and $\lambda_{\max}(Q)$ for its smallest and largest
  eigenvalues, and
  $\kappa(Q)=\lambda_{\max}(Q)/\lambda_{\min}(Q)$ for its spectral condition
  number.
  \item A matrix $A\in\R^{m\times d}$ has \emph{full row rank} when
  $\operatorname{rank}(A)=m$. Its singular values are the nonnegative square
  roots of the eigenvalues of $A\trans A$.
  \item For a descent direction $p_k$, Armijo backtracking starts from a trial
  step $\alpha>0$ and repeatedly replaces $\alpha$ by $\rho\alpha$, with
  $0<\rho<1$, until
  \[
    f(x_k+\alpha p_k)
    \leq f(x_k)+c\alpha\nabla f(x_k)\trans p_k,
    \qquad 0<c<1.
  \]
  \item A centered finite-difference approximation of the $j$th gradient
  component is
  \[
    \partial_j f(x)\approx
    \frac{f(x+he_j)-f(x-he_j)}{2h},
  \]
  where $e_j$ is the $j$th coordinate vector and $h>0$ is small.
  \item A first-order perturbation keeps only terms linear in a small change.
  In particular,
  \[
    (\Sigma+E)^{-1}
    \approx\Sigma^{-1}-\Sigma^{-1}E\Sigma^{-1}
  \]
  when $E$ is small and $\Sigma$ is invertible.
  \item \emph{In-sample} data are used to estimate or select a model;
  \emph{out-of-sample} data are held back and used only for final evaluation.
  A hat, as in $\widehat\mu$ or $\widehat\Sigma$, denotes an estimate.
\end{itemize}
\end{conceptbox}

% -----------------------------------------------------------------------------
\section{Part 1 -- One-Dimensional Unconstrained Optimization}

\paragraph{Exercise 1 -- Critical points are only candidates.}
For each function below, find every critical point, classify it as a local
minimum, local maximum, or neither, and determine whether a global optimizer
exists on the stated domain:
\[
f_1(x)=x^4-4x^2+3,\quad x\in\R,
\]
\[
f_2(x)=x^3-3x+1,\quad x\in[-2,2],
\qquad
f_3(x)=\log(1+e^x)-\frac34x,\quad x\in\R.
\]
\begin{enumerate}
\item Use first- and second-order conditions where they apply.
\item For $f_2$, include the boundary points in the global comparison.
\item Prove that $f_3$ is strictly convex and compute its unique minimizer.
\item Sketch the three functions and mark all candidates.
\end{enumerate}
\paragraph{Exercise 2 -- A parameter changes the number of local optima.}
Let
\[
f_a(x)=\frac14x^4-\frac a2x^2+x,
\qquad a>0.
\]
\begin{enumerate}
\item Show that stationary points solve $x^3-ax+1=0$.
\item Study the derivative of this cubic and determine for which values of
$a$ the objective has one stationary point and for which values it has three.
\item In the three-point regime, classify the stationary points without
solving the cubic explicitly.
\item Without computing decimal approximations, explain how the sign pattern
of $f_a'$ orders the local extrema in the three-point regime.
\item Explain why $f_a'(x)=0$ alone cannot identify the global solution and
give a rigorous comparison strategy.
\end{enumerate}

\paragraph{Exercise 3 -- Newton's method as a local model.}
Consider
\[
f(x)=e^x-3x+\frac14x^2.
\]
\begin{enumerate}
\item Prove that $f$ has a unique global minimizer.
\item Write Newton's iteration for solving $f'(x)=0$.
\item Express the Newton direction as the minimizer of the local quadratic
model of $f$ at $x_k$.
\item Determine conditions on $x_k$ under which the full Newton step is a
descent direction.
\item Explain how Armijo backtracking modifies the update and why it improves
global robustness. Numerical trajectories belong to Exercise 10.
\end{enumerate}

% -----------------------------------------------------------------------------
\section{Part 2 -- Multidimensional Unconstrained Optimization}

\paragraph{Exercise 4 -- Quadratic optimization and conditioning.}
Let
\[
f(x)=\frac12x\trans Qx-c\trans x,\qquad
Q=Q\trans\succ0,\qquad x,c\in\R^d.
\]
\begin{enumerate}
\item Compute $\nabla f$ and $\nabla^2f$.
\item Proceed in three steps:
\begin{enumerate}[label=(\alph*)]
  \item use $Q\succ0$ to prove that $f$ is strictly convex;
  \item write the stationarity equation as the linear system
  $Qx^\star=c$;
  \item explain why $Q\succ0$ implies that $Q$ is invertible and conclude that
  this system has exactly one solution.
\end{enumerate}
The notation $x^\star=Q^{-1}c$ may be used for the theoretical formula, but a
numerical implementation should solve $Qx=c$ rather than form $Q^{-1}$.
\item Complete the square to express $f(x)-f(x^\star)$ as a quadratic form.
\item In an orthonormal eigenbasis of $Q$, interpret the eigenvalues as
principal curvatures.
\item Let $e_k=x_k-x^\star$ and consider gradient descent with a constant step
$\alpha$. Detail the role of the condition number as follows:
\begin{enumerate}[label=(\alph*)]
  \item show in an orthonormal eigenbasis of $Q$ that each error coordinate
  satisfies
  \[
    (e_{k+1})_i=(1-\alpha\lambda_i)(e_k)_i;
  \]
  \item explain why a large ratio
  $\kappa(Q)=\lambda_{\max}(Q)/\lambda_{\min}(Q)$ produces elongated level
  sets and very different contraction factors across directions;
  \item recall the stability condition $0<\alpha<2/\lambda_{\max}(Q)$ and
  compare qualitatively the cases $\kappa(Q)\approx1$ and
  $\kappa(Q)\gg1$.
\end{enumerate}
\end{enumerate}

\paragraph{Exercise 5 -- Least squares, ridge regression, and conditioning.}
Let $A\in\R^{n\times d}$ and $b\in\R^n$.
\begin{enumerate}
\item Minimize $\frac12\norm{Ax-b}_2^2$ by deriving the normal equations.
\item Give a necessary and sufficient condition for uniqueness. Interpret the
residual geometrically when the normal equations hold.
\item For $\lambda>0$, solve
\[
\min_x\ \frac12\norm{Ax-b}_2^2+\frac\lambda2\norm{x}_2^2.
\]
\item Here, \emph{rank deficient} means that $A$ does not have full column
rank:
\[
  \operatorname{rank}(A)<d.
\]
Equivalently, the columns of $A$ are linearly dependent and $A\trans A$ is
singular. Show that ridge regression nevertheless has a unique solution by
proving that $A\trans A+\lambda I\succ0$ for every $\lambda>0$.
\item Use the singular values of $A$ to explain how regularization changes the
condition number and creates a bias--stability trade-off.
\end{enumerate}

\paragraph{Exercise 6 -- Second-order conditions in several dimensions.}
For each function, find all stationary points and classify them:
\[
f_1(x,y)=x^2+xy+y^2-4x+2y,
\]
\[
f_2(x,y)=x^4+y^4-2x^2-4y^2,
\qquad
f_3(x,y)=x^2-y^2+2xy.
\]
\begin{enumerate}
\item Compute each Hessian and use its eigenvalues or principal minors.
\item For degenerate candidates of $f_2$, inspect the objective along useful
directions when the Hessian test is inconclusive.
\item Identify which local minima are also global and justify the answer.
\item Plot contours for the three objectives on a domain of your choice.
\end{enumerate}

% -----------------------------------------------------------------------------
\section{Part 3 -- Multidimensional Constrained Optimization}

\paragraph{Exercise 7 -- Projection onto an affine constraint.}
Let $z,a\in\R^d$, with $a\neq0$, and
$C=\{x\in\R^d:a\trans x=b\}$.
\begin{enumerate}
\item Formulate $\proj_C(z)$ as a constrained quadratic problem.
\item Solve it with a Lagrange multiplier.
\item Obtain a closed formula for $\proj_C(z)$ and verify that it returns $z$
when $z\in C$.
\item Verify geometrically that $z-\proj_C(z)$ is normal to $C$.
\item Extend the derivation to $C=\{x\in\R^d:Ax=b\}$, where
$A\in\R^{m\times d}$. The matrix $A$ has \emph{full row rank} when
\[
  \operatorname{rank}(A)=m,
\]
meaning that its $m$ rows are linearly independent. Explain why this implies
that $AA\trans$ is invertible, then derive the projection formula.
\end{enumerate}

\paragraph{Exercise 8 -- Equality constraints, duality, and a saddle point.}
Consider the generic equality-constrained quadratic problem
\[
\min_{x\in\R^d}\ \frac12x\trans Qx+c\trans x
\quad\text{subject to}\quad Ax=b,
\]
where $Q\in\R^{d\times d}$ is symmetric positive definite and
$A\in\R^{m\times d}$ has full row rank.
\begin{enumerate}
\item Write the Lagrangian and the KKT linear system.
\item Eliminate $x$ and derive formulas for $\lambda^\star$ and $x^\star$.
\item Compute the dual function
\[
q(\lambda)=\inf_x\mathcal L(x,\lambda)
\]
explicitly and show directly that it is concave.
\item Prove weak duality for every feasible $x$ and every multiplier
$\lambda$.
\item Verify strong duality and the saddle inequalities symbolically.
\item With the convention
$\mathcal L(x,\lambda)=f(x)+\lambda\trans(Ax-b)$, show that the first-order
sensitivity of the optimal value to $b$ is $-\lambda^\star$.
\end{enumerate}

\paragraph{Exercise 9 -- Projection onto a half-space with KKT.}
Let $z,a\in\R^d$, $a\neq0$, and solve
\[
\min_x\ \frac12\norm{x-z}_2^2
\quad\text{subject to}\quad a\trans x\leq b.
\]
\begin{enumerate}
\item Draw the geometry in dimension two and identify the unconstrained
minimizer.
\item Write every KKT condition, including dual feasibility and
complementary slackness.
\item Treat separately the cases $a\trans z\leq b$ and $a\trans z>b$.
Derive the optimizer and multiplier in each case.
\item Explain why the KKT conditions are sufficient here.
\item Show that the projection map can be written using the positive part
$(u)_+=\max(u,0)$ and interpret the active-set change.
\end{enumerate}

% -----------------------------------------------------------------------------
\section{Part 4 -- Python Algorithms}

\paragraph{Exercise 10 -- Implement gradient descent from scratch.}
Implement
\[
x_{k+1}=x_k-\alpha_k\nabla f(x_k)
\]
without calling an optimization routine. For the quadratic test, set a random
seed and construct $Q=M\trans M+0.1I$ and $c$ with $d=20$; retain the exact
solution from Exercise 4 as a benchmark. Also test the Rosenbrock function
\[
r(x,y)=100(y-x^2)^2+(1-x)^2.
\]
\begin{enumerate}
\item Write a reusable function accepting \texttt{f}, \texttt{grad}, an
initial point, a step-size rule, a tolerance, and a maximum iteration count.
\item Return the entire history of iterates, objective values, gradient
norms, and accepted step sizes.
\item For the quadratic, compare constant steps $1/\lambda_{\max}(Q)$,
$2/(\lambda_{\min}(Q)+\lambda_{\max}(Q))$, and one unstable step satisfying
$\alpha>2/\lambda_{\max}(Q)$.
\item Add Armijo backtracking with $c=10^{-4}$ and $\rho=1/2$, then test the
code on Rosenbrock from $(-1.2,1)\trans$.
\item Plot objective gap and gradient norm on logarithmic axes. Plot the
iterates over objective contours.
\item Check the analytical gradient with centered finite differences before
running the optimizer.
\end{enumerate}

\begin{deliverablebox}{Algorithmic deliverable}
Submit the implementation, two convergence tables, and four diagnostic plots.
State the stopping rule explicitly. A plot alone is not evidence of
convergence: compare with the known solution whenever one is available.
\end{deliverablebox}

\paragraph{Exercise 11 -- Implement primal--dual gradient descent--ascent.}
Use the equality-constrained quadratic from Exercise 8. Generate a reproducible
instance with $d=20$, $m=3$, $Q=M\trans M+0.1I$, and a full-row-rank matrix
$A$. Draw $c$ and a reference point $x^{\mathrm{ref}}$, then set
$b=Ax^{\mathrm{ref}}$ so feasibility is guaranteed. Its Lagrangian is
\[
\mathcal L(x,\lambda)
=\frac12x\trans Qx+c\trans x+\lambda(Ax-b).
\]
Implement the simultaneous iteration
\[
x_{k+1}=x_k-\alpha\bigl(Qx_k+c+A\trans\lambda_k\bigr),
\qquad
\lambda_{k+1}=\lambda_k+\beta(Ax_k-b).
\]
\begin{enumerate}
\item Start from $x_0=0$, $\lambda_0=0$ and test several pairs
$(\alpha,\beta)$.
\item At every iteration record the primal error, feasibility, and stationarity
residuals
\[
\norm{x_k-x^\star},\qquad
\norm{Ax_k-b},\qquad
\norm{Qx_k+c+A\trans\lambda_k},
\]
as well as the Lagrangian value.
\item Find one stable and one oscillatory parameter choice. Plot the
trajectories and explain the difference.
\item Implement the augmented Lagrangian
\[
\mathcal L_\rho(x,\lambda)=\mathcal L(x,\lambda)
+\frac\rho2\norm{Ax-b}_2^2
\]
and its descent--ascent updates. Study $\rho=0.1,1,10$.
\item Compare all numerical solutions with the exact KKT solution from
Exercise 8.
\item Optional: define $z=(x,\lambda)$ and the saddle operator
$F(z)=(\nabla_x\mathcal L,-\nabla_\lambda\mathcal L)$. Implement the
extragradient update
\[
  \widetilde z_k=z_k-\eta F(z_k),
  \qquad z_{k+1}=z_k-\eta F(\widetilde z_k),
\]
then compare its stability with simultaneous descent--ascent.
\end{enumerate}

% -----------------------------------------------------------------------------
\section{Part 5 -- Portfolio Optimization}

Let $R\in\R^d$ be the random vector of annual asset returns, with
mean $\mu=\E[R]$ and covariance matrix $\Sigma=\Cov(R)$. Unless otherwise
stated, portfolio weights satisfy $\mathbf1\trans w=1$. For risk-free rate
$r_f$, define
\[
R(w)=\mu\trans w,
\qquad \sigma(w)=\sqrt{w\trans\Sigma w},
\qquad \operatorname{SR}(w)=\frac{R(w)-r_f}{\sigma(w)}.
\]

\paragraph{Exercise 12 -- Risk, return, and the Sharpe ratio.}
Let $\mu\in\R^d$, $\Sigma\succ0$, and $r_f\in\R$.
\begin{enumerate}
\item Prove the formulas for portfolio return and variance from the random
return vector $R$.
\item Derive the global minimum-variance portfolio allowing short sales:
\[
w_{\rm GMV}=\frac{\Sigma^{-1}\mathbf1}
{\mathbf1\trans\Sigma^{-1}\mathbf1}.
\]
\item Show that the tangency direction is proportional to
$\Sigma^{-1}(\mu-r_f\mathbf1)$ and explain how full investment fixes its
normalization when the denominator is nonzero.
\item Compare the objectives answered by equal weight, GMV, and tangency
portfolios. State when each is a sensible benchmark.
\item Explain why long-only constraints can invalidate the closed formulas
and identify the KKT conditions required instead.
\end{enumerate}

\paragraph{Exercise 13 -- Derive the efficient frontier.}
Consider
\[
\min_w\ \frac12w\trans\Sigma w
\quad\text{subject to}\quad
\mathbf1\trans w=1,\qquad \mu\trans w=m
\]
for a feasible target return $m$.
\begin{enumerate}
\item Derive the KKT system for a fixed target $m$.
\item With
\[
C=\begin{pmatrix}\mathbf1\trans\\\mu\trans\end{pmatrix},
\qquad d_m=\begin{pmatrix}1\\m\end{pmatrix},
\]
derive
$w(m)=\Sigma^{-1}C\trans(C\Sigma^{-1}C\trans)^{-1}d_m$.
\item Show that the optimal variance is a quadratic function of $m$ and
deduce the shape of the frontier in mean--standard-deviation space.
\item Explain qualitatively why long-only constraints and position caps create
kinks and active-set changes.
\item State the feasibility and positive-semidefiniteness checks a numerical
solver should perform before tracing a frontier.
\end{enumerate}

\paragraph{Exercise 14 -- Estimation error and robust weights.}
Suppose an estimate $(\widehat\mu,\widehat\Sigma)$ replaces the population
parameters, with $\widehat\Sigma\succ0$. Define the unconstrained tangency
direction
\[
q(\widehat\mu,\widehat\Sigma)
=\widehat\Sigma^{-1}(\widehat\mu-r_f\mathbf1).
\]
\begin{enumerate}
\item Study the sensitivity of $q$ to mean-estimation error. Throughout this
question, keep $\widehat\Sigma$ and $r_f$ fixed. Let $h\in\mathbb R^n$ denote
an estimation-error direction and replace $\widehat\mu$ by
$\widehat\mu+\varepsilon h$.
  \begin{enumerate}[label=(\alph*)]
  \item Write $q(\widehat\mu+\varepsilon h,\widehat\Sigma)$ and isolate the
  first-order change
  $q(\widehat\mu+\varepsilon h,\widehat\Sigma)
  -q(\widehat\mu,\widehat\Sigma)$.
  \item Deduce the directional derivative
  $D_{\mu}q(\widehat\mu,\widehat\Sigma)[h]$ and the Jacobian of $q$ with
  respect to $\widehat\mu$.
  \item Use an eigendecomposition
  $\widehat\Sigma=U\operatorname{diag}(\lambda_1,\ldots,\lambda_n)U\trans$.
  Decompose $h$ in the orthonormal eigenbasis and determine how the component
  along an eigenvector associated with $\lambda_i$ is transformed.
  \item Give an upper bound for the relative or absolute change in $q$ in
  terms of $\lambda_{\min}(\widehat\Sigma)$ and $\lVert h\rVert_2$. Explain
  why a mean-estimation error aligned with a small-eigenvalue direction can
  produce a large change in the portfolio direction.
  \item Optional: if the fully invested weights are
  $w=q/(\mathbf1\trans q)$, explain why the normalization may create an
  additional instability when $\mathbf1\trans q$ is close to zero.
  \end{enumerate}
\item Explain why GMV weights are insensitive to mean estimation but remain
sensitive to covariance estimation.
\item Define the ridge-regularized covariance
\[
\widehat\Sigma_\delta=\widehat\Sigma+\delta I,
\qquad \delta\geq0.
\]
Show how its eigenvalues and condition number depend on $\delta$.
\item Compare in-sample predicted volatility with out-of-sample realized
volatility. Explain why optimizing estimated expected returns is particularly
unstable.
\item Propose a transparent validation rule for selecting $\delta$ without
looking at the final test sample. Reserve all numerical comparisons for the
separate portfolio project.
\end{enumerate}

% -----------------------------------------------------------------------------
\begin{deliverablebox}{Separate Portfolio Optimization Project}
The integrated project is in the separate Portfolio Optimization Project
folder. Use its detailed statement and \path{portfolio_project_skeleton.py},
which contains the staged questions, validation checks, expected outputs, and
explicit \texttt{TODO} blocks.
\end{deliverablebox}

\end{document}
